% Guia do professor — Capítulo 20: Previsão Numérica do Tempo
\section{Capítulo 20 --- Previsão Numérica do Tempo}

\subsection*{Visão geral e ritmo}

O capítulo-síntese computacional: as 7 equações do curso viram
código, e as quatro peças (discretizar, parametrizar, inicializar,
quantificar) organizam tudo. Para físicos é território natal —
métodos numéricos, caos, inferência — aplicado ao maior problema de
valor inicial do mundo. \textbf{Ritmo: 2 aulas de 2\,h.}

\begin{itemize}
  \item \textbf{Aula 1} — Richardson$\to$ENIAC (10 min de história
        que valem ouro), grade, esquemas e CFL (Caso~1 ao vivo),
        parametrizações como física reduzida.
  \item \textbf{Aula 2} — Lorenz e o horizonte (Caso~2), assimilação
        (Caso~3), ensembles e probabilidade (Caso~4), a revolução
        silenciosa.
\end{itemize}

\subsection*{Objetivos de aprendizagem}

\begin{enumerate}
  \item listar as equações primitivas e reconhecê-las como o curso
        inteiro;
  \item aplicar CFL e prever os erros característicos de cada
        esquema;
  \item explicar o horizonte de ~2 semanas via Lyapunov;
  \item usar o ganho ótimo da assimilação e explicar por que a
        análise supera modelo e observação;
  \item interpretar previsões probabilísticas de ensemble.
\end{enumerate}

\subsection*{Pré-requisitos a reativar}

As 7 equações espalhadas (Caps.~1, 3, 4, 10); advecção numérica
(Cap.~3, N3!); tudo que virou parametrização (Caps.~2, 7, 14, 17,
18); mínimos quadrados/propagação de erros (laboratório de física).

\subsection*{Armadilhas conceituais frequentes}

\begin{itemize}
  \item \textbf{``Mais computador $=$ mais previsão''}: o horizonte é
        logarítmico no erro inicial (T3) — caos é estrutural, não
        falta de hardware.
  \item \textbf{``O modelo de 10\,km resolve 10\,km''}: resolução
        efetiva $\sim7\Delta x$; cumulus continuam parametrizados.
  \item \textbf{Assimilação como ``usar os dados''}: é média
        ponderada por precisões — jogar o modelo fora e usar só
        observações seria PIOR (T4).
  \item \textbf{``70\% de chuva'' como enfeite}: é a fração literal
        do ensemble — probabilidade é a saída primária, não o
        determinístico.
  \item \textbf{Caos como ruído}: o atrator é determinístico e
        estruturado — imprevisível no ponto, previsível na
        estatística (ponte para o Cap.~21).
\end{itemize}

\subsection*{Demonstração de baixo custo}

Rodar o Caso~2 ao vivo mudando o 6º decimal — o momento
``uau'' garantido do semestre. Comparar a previsão de 5 dias do
ECMWF/GFS de ontem com a carta de hoje; abrir o meteograma de
ensemble do ECMWF para a cidade (espaguete real).

\subsection*{Problemas sugeridos do livro (exercícios do Cap.~20)}

Aquecimento: CFL para grades/ventos dados; contas de pontos de grade
e memória. Aprofundamento: von Neumann de outros esquemas; ganho
ótimo com números. Síntese: ``por que o clima de 2100 é mais
previsível que o tempo de daqui a 20 dias?''.
\emph{Confira a numeração na sua cópia (v1.02b).}

\subsection*{Uso do notebook \texttt{cap20\_previsao.ipynb}}

\begin{center}
\begin{tabular}{lll}
\hline
Caso & Conteúdo & Momento sugerido \\
\hline
1 & upwind/Lax--Wendroff/leapfrog $+$ CFL & Aula 1 \\
2 & Lorenz 63: borboleta e gêmeos & Aula 2, abertura \\
3 & ciclo de assimilação no Lorenz & Aula 2 \\
4 & ensemble: espaguete $\to$ probabilidade & Aula 2, fecho \\
\hline
\end{tabular}
\end{center}

O arco dos Casos 2$\to$3$\to$4 é a história intelectual da previsão
moderna em 30 minutos de aula.

\subsection*{Exercícios complementares e soluções}

Nas notas: T1--T5 e N1--N4. Soluções em \texttt{cap20\_solucoes.ipynb}
(\emph{não distribuir}): T1, T2 e T4 com \texttt{sympy} (von Neumann,
equação modificada e ganho ótimo simbólicos); N4 é a pérola
(dispersão prevê o erro). Pares: T1/T2$\leftrightarrow$N1,
T3$\leftrightarrow$N2, T4$\leftrightarrow$N3.

\subsection*{Conexões com o resto do curso}

TODAS: as equações (Caps.~1, 3, 4, 10), as parametrizações
(Caps.~2, 7, 14, 17, 18), as observações (Caps.~8--9), os fenômenos
que o modelo deve acertar (Caps.~11--16) e a fronteira com o clima
(Cap.~21). Este capítulo é o ponto de fuga do curso.
